\documentclass[11pt]{article}
\usepackage[hmargin=1.05in,vmargin=0.8in]{geometry}
\usepackage{amsmath,amssymb}
\usepackage[hidelinks]{hyperref}

\title{A counterexample to Conjecture 00000008827}
\author{Feiyu (Independent researcher)}
\date{}

\begin{document}
\maketitle

\begin{abstract}
For one empirical observation, an affine loss, and compact convex
support, the $1$-Wasserstein robust dual has the same value at radii
$1$ and $2$. Both values exceed the empirical risk. Hence the dual
cannot equal the empirical risk plus the radius times a fixed penalty.
\end{abstract}

\section*{The regularization claim}
The final clause of the \href{https://github.com/The-Last-Math-Competition/The-Last-Math-Competition/blob/main/conjectures/00000008827.md}{English and Chinese statements}
asserts exact equality of the Wasserstein robust dual and a regularized
empirical risk with coefficient equal to the ball radius. With the
loss, data, and decision held fixed, this coefficient claim requires
$D(r)=\widehat R+rP$ for one penalty value $P$ and all $r>0$.

\section*{An exact dual calculation}
Let the support be $[-2,1]$ with distance $|z-z'|$, take the single
observation $0$, and use the nonnegative affine loss $\ell(z)=z+2$.
Thus $\widehat R=\ell(0)=2$. Every transport plan from the empirical
point mass $\delta_0$ has first coordinate $0$ almost surely, so its
cost is the expected value of $|z|$. The transport dual for radius $r$
is therefore
\[
 D(r)=\inf_{\lambda\ge 0}
   \left\{\lambda r+\sup_{z\in[-2,1]}\bigl(z+2-\lambda|z|\bigr)\right\}.
\]
This is the standard one-sample Wasserstein dual, also obtained by
dualizing the expected-cost cap with multiplier $\lambda$.

For $r\ge 1$, evaluating the inner expression at $z=1$ gives the lower
bound $3+\lambda(r-1)\ge 3$. Taking $\lambda=0$ gives exactly $3$.
The worst-case distribution $\delta_1$ also attains loss $3$ at cost
$1$. In particular, $D(1)=D(2)=3$ while $\widehat R=2$. The proposed
identity would require both $P=1$ and $P=1/2$.
Both radii lie below the support diameter of $3$.

\section*{Lean verification}
The Lean project represents the interval as a subtype of real
numbers and defines the transport dual with its infimum over
nonnegative multipliers and its supremum over the support.
The final theorem refutes the fixed-penalty formula using the two
dual values. Build with \texttt{lake build} in \texttt{lean/} using
Lean and Mathlib v4.33.1. The axiom dependencies are
\texttt{propext}, \texttt{Classical.choice}, and \texttt{Quot.sound}.

\smallskip
\noindent
Reference. P. Mohajerin Esfahani and D. Kuhn,
\href{https://doi.org/10.1007/s10107-017-1172-1}{Data-driven distributionally robust optimization using the Wasserstein metric},
\emph{Mathematical Programming} 171 (2018), 115--166.
\end{document}
